% Generated draft: requires mathematical and source review.
\documentclass[11pt]{article}
\usepackage[margin=25mm]{geometry}
\usepackage{fontspec}
\setmainfont{Latin Modern Roman}
\usepackage{amsmath,amssymb,mathtools,unicode-math}
\setmathfont{Latin Modern Math}
\usepackage{enumitem,needspace,xurl,hyperref}
\setlength{\parindent}{0pt}
\setlength{\parskip}{5pt}
\setlength{\emergencystretch}{3em}
\newcommand{\N}{\mathbb{N}}
\newcommand{\Z}{\mathbb{Z}}
\newcommand{\Q}{\mathbb{Q}}
\newcommand{\R}{\mathbb{R}}
\newcommand{\C}{\mathbb{C}}
\newcommand{\F}{\mathbb{F}}
\newcommand{\A}{\mathbb{A}}
\newcommand{\PP}{\mathbb{P}}
\newcommand{\E}{\mathbb{E}}
\newcommand{\eps}{\varepsilon}
\newcommand{\dd}{\,\mathrm{d}}
\begin{document}
\section*{UnsolvedMath ID: 801148}
\textbf{Source problem code:} catalogue entry
\subsection{Definitions and mathematical statement}
The supplied record concerns the Korteweg--de Vries (KdV) system and displays the expression
\[
V(t)=\frac{C(t)}{m(t)}.
\]
This quotient is meaningful only where its denominator is nonzero. The record does not define \(C(t)\), \(m(t)\), their domains or regularity, or the construction by which \(V\) generates a potential.

\textbf{Problem as supplied.} Determine what is special about the potential generated by \(V(t)=C(t)/m(t)\).

\textbf{Supplied background and qualifications.} The source authors report that, for the KdV system, an unspecified preceding task was already carried out in reference [veselov], although they did not find detailed formulas. They describe [moser81] as relating stationary KdV solutions to solutions of Neumann systems: the potential energy of a Neumann system, evaluated along a trajectory and suitably scaled, yields a solution of the KdV system. They further state that the KdV solutions obtained this way are related to potentials in the Hill--Schr\"odinger equation whose spectrum has finitely many bands. The supplied text also refers to [DKN, p.~219, equation~(2.162), pp.~247--248]. These are attributed background assertions, not conclusions established here.

\textbf{Unresolved formulation.} The supplied text does not specify the preceding task, the relevant Neumann system, the scaling, the KdV equation or stationary reduction, or the relation between \(V\) and a Hill--Schr\"odinger potential. It also does not identify the property requested by ``what is special.'' Consequently, the displayed question cannot be turned into a uniquely specified mathematical assertion from this record alone. In particular, no spectral property or boundary conditions for the generated potential are stipulated here.
\subsection{Short English statement}
Problem 1 asks what distinguishes the potential generated by \(V(t)=C(t)/m(t)\), against a background connecting Neumann systems, stationary KdV solutions, and Hill--Schr\"odinger potentials with finitely many spectral bands. The functions and the generation procedure are missing, so the intended characterization remains unresolved.
\subsection{Sources}
\begin{itemize}
\item Local registry: \texttt{data/unsolvedmath/problems.json}, permanent UnsolvedMath ID \(801148\), exact source code \texttt{catalogue entry}; supplied title: \emph{Nijenhuis Geometry and Integrable Systems --- Problem 1}. Attribution: Bolsinov, Konyaev, Matveev; proposed year: 2024. The registry labels the problem open; that label is supplied metadata and has not been independently verified.
\item Bolsinov, Konyaev, Matveev, \emph{Research Problems on Nijenhuis Geometry and Integrable Systems} (2024), Problem 1, \url{https://arxiv.org/abs/2410.04276}. The supplied provenance records access on 2026-10-16 and extraction by \texttt{arxiv\_latex}; these are source-record provenance claims, not a claim of access or verification for this draft.
\item Reference label [veselov], cited in the supplied excerpt for prior work concerning the KdV system. The source authors qualify their report by saying that they did not find detailed formulas. No bibliographic entry for this label is supplied.
\item Reference label [moser81], described in the supplied excerpt as a booklet relating stationary KdV solutions and Neumann systems. No bibliographic entry for this label is supplied.
\item Reference label [DKN], p.~219, equation~(2.162), pp.~247--248. No bibliographic entry for this label is supplied.
\end{itemize}
\end{document}
